\documentclass[12pt,a4paper]{article}

\usepackage[utf8]{inputenc}
\usepackage[T1]{fontenc}
\usepackage[spanish]{babel}
\usepackage[top=2.5cm,bottom=2.5cm,left=2.5cm,right=2.5cm,headheight=15pt]{geometry}
\usepackage{setspace}
\onehalfspacing
\usepackage{xcolor}
\definecolor{darkblue}{RGB}{0,0,139}
\usepackage{fancyhdr}
\pagestyle{fancy}
\fancyhf{}
\fancyhead[L]{\small\textit{Dota 2 Pro Matches --- An\'alisis Exploratorio de Datos}}
\fancyhead[R]{\small T\'opicos en Ciencias de Datos}
\fancyfoot[C]{\thepage}
\renewcommand{\headrulewidth}{0.4pt}
\usepackage{titlesec}
\titleformat{\section}{\large\bfseries}{}{0em}{}[\titlerule]
\titleformat{\subsection}{\normalsize\bfseries}{}{0em}{}
\titleformat{\subsubsection}{\normalsize\bfseries\itshape}{}{0em}{}
\titlespacing*{\section}{0pt}{14pt}{6pt}
\titlespacing*{\subsection}{0pt}{10pt}{4pt}
\usepackage{booktabs}
\usepackage{array}
\usepackage{tabularx}
\usepackage{longtable}
\usepackage{enumitem}
\setlist[itemize]{topsep=3pt,itemsep=2pt,parsep=0pt}
\usepackage{graphicx}
\usepackage{caption}
\captionsetup{font=small,labelfont=bf,skip=4pt}
\usepackage{amsmath}
\usepackage{amssymb}
\usepackage{parskip}
\setlength{\parskip}{6pt}
\usepackage{mdframed}
\usepackage{tocloft}
\usepackage{hyperref}
\hypersetup{colorlinks=true,linkcolor=darkblue,urlcolor=darkblue,citecolor=darkblue}
\renewcommand{\cfttoctitlefont}{\large\bfseries}
\setlength{\cftbeforesecskip}{4pt}

% Caja para placeholders de imágenes
% Para reemplazar por la imagen real, sustituir el bloque \imagenplaceholder{...}{...}
% por: \includegraphics[width=0.95\linewidth]{graficoN.png}
\newcommand{\imagenplaceholder}[2]{%
  \begin{center}
    \fbox{%
      \begin{minipage}{0.85\linewidth}
        \centering
        \vspace{1.5cm}
        {\large\bfseries #1}\\[0.4em]
        {\small\itshape #2}
        \vspace{1.5cm}
      \end{minipage}%
    }
  \end{center}
}

\begin{document}

%% PORTADA
\begin{titlepage}
  \thispagestyle{empty}
  \begin{flushleft}
    Universidad La Salle\\
  \end{flushleft}
  \vspace*{\fill}
  \begin{center}
    {\Large\bfseries INFORME DE AN\'ALISIS EXPLORATORIO DE DATOS}\\[1.2em]
    {\Large Dota 2 Pro Matches}\\[1.5em]
    {\normalsize Inteligencia de Negocios y Miner\'ia de Datos}
  \end{center}
  \vspace*{\fill}
  \begin{flushright}
    \begin{tabular}{ll}
      Integrante: & Piero Adrian Delgado Chipana \\[0.4em]
      Profesor:   & Ana Maria Cuadros Valdivia   \\[0.4em]
      \multicolumn{2}{l}{Arequipa --- Per\'u}
    \end{tabular}
  \end{flushright}
\end{titlepage}

%% INDICE
\thispagestyle{fancy}
\renewcommand{\contentsname}{\'Indice de Contenidos}
\tableofcontents
\newpage

%% ============================================================
\section{Definici\'on del Problema}
%% ============================================================

Dota~2 es un videojuego de estrategia en l\'inea (MOBA) de \textbf{5 jugadores
contra 5}, donde cada jugador controla un \textbf{h\'eroe} con habilidades
\'unicas. Un equipo se denomina \textbf{Radiant} y el otro \textbf{Dire}; la
partida termina cuando un equipo destruye el \textit{Ancient} rival.

Entre 2019 y mediados de 2021, Dota~2 atraves\'o su \textbf{periodo de mayor
auge profesional}: el juego cont\'o con el mayor n\'umero de equipos activos,
torneos internacionales y audiencia en toda su historia. En esos a\~nos el
juego ten\'ia una versi\'on del mapa, un pool de h\'eroes y un sistema de \'items
cuyo meta competitivo fue ampliamente estudiado. Desde entonces, Dota~2 recibi\'o
cambios radicales en el dise\~no del mapa, nuevos h\'eroes y reworks completos de
mec\'anicas.

La \textbf{problem\'atica principal} es que los jugadores y analistas que hoy
buscan referencias del juego profesional de esa \'epoca ---para estudiar estrategias,
entender el meta o prepararse para competir--- no disponen de una herramienta
que les permita explorar esas partidas de forma estructurada.

%% ============================================================
\section{Objetivos}
%% ============================================================

\begin{itemize}
  \item \textbf{Predictivo:} demostrar que el equipo Radiant presenta ventajas
        estructurales sobre Dire incluso antes de que la partida comience,
        cuantificables a trav\'es del rendimiento hist\'orico acumulado de cada
        equipo (GPM, XPM, KDA, win rate, etc.) y verificables con la variable
        objetivo \texttt{radiant\_win}.
  \item \textbf{Exploratorio y manejo de errores:} identificar y gestionar
        las inconsistencias del dataset: partidas no completadas por abandono,
        partidas de torneo jugadas simult\'aneamente con distinto \texttt{match\_id}
        pero mismo resultado y rendimiento (y m\'as h\'eroes registrados), y
        partidas sin h\'eroes o con h\'eroes incompletos que deben excluirse del
        an\'alisis t\'actico pero no del an\'alisis de rendimiento general.
\end{itemize}

%% ============================================================
\section{Hip\'otesis Iniciales}
%% ============================================================

\subsection{Motivaci\'on}

Las hip\'otesis surgieron a partir de la observaci\'on de que los equipos
profesionales desarrollan ventajas medibles en distintas dimensiones del juego:
eficiencia econ\'omica, rendimiento de combate, control de visi\'on y objetivos
neutrales. Al explorar las 66 columnas de estad\'isticas del dataset se
identificaron variables con potencial predictivo directo sobre el resultado de la
partida (\texttt{radiant\_win}). Adicionalmente, el sesgo estructural de mapa de
Radiant (52.1\,\% de victorias) y la ley de los grandes n\'umeros sobre ventanas
hist\'oricas grandes motivaron formular hip\'otesis verificables.

\subsection{Hip\'otesis en Forma de Pregunta}

\textbf{Hip\'otesis 1:}\\
\textit{¿Los equipos con mayor \texttt{gold\_per\_min\_avg} y \texttt{xp\_per\_min\_avg}
tienen mayor probabilidad de ganar, siendo la eficiencia econ\'omica el principal
predictor de victoria en Dota~2 profesional?}

\textbf{Hip\'otesis 2:}\\
\textit{¿Los equipos con mejor KDA hist\'orico (\texttt{kda\_avg} $>$ 3.5) tienen una
tasa de victoria (\texttt{win\_pct}) significativamente mayor que equipos con
KDA $<$ 2.0?}

\textbf{Hip\'otesis 3:}\\
\textit{¿El control de visi\'on --- medido por la cantidad de wards colocadas y
destruidas --- diferencia estad\'isticamente a los equipos ganadores de los perdedores?}

\textbf{Hip\'otesis 4:}\\
\textit{¿Matar a Roshan (\texttt{roshan\_kills\_avg} alto) est\'a asociado a una
mayor probabilidad de ganar, especialmente en partidas disputadas de larga duraci\'on?}

\textbf{Hip\'otesis 5:}\\
\textit{¿El control de torres (\texttt{tower\_kills\_avg}) es m\'as predictivo del
resultado que el control de kills de h\'eroes (\texttt{hero\_kills\_avg})?}

\subsection{Plan de An\'alisis}

Para investigar cada hip\'otesis se sigui\'o el siguiente proceso:

\begin{enumerate}
  \item \textbf{Carga y limpieza:} lectura del CSV, detecci\'on de partidas sin
        historial ($\sim$6\,\%), imputaci\'on con la media columnar.
  \item \textbf{An\'alisis descriptivo:} c\'alculo de medidas de tendencia central
        y dispersi\'on para las 66 features principales.
  \item \textbf{Distribuci\'on de resultados:} conteo de victorias Radiant vs.\ Dire
        por a\~no y global.
  \item \textbf{An\'alisis de correlaci\'on:} relaci\'on entre GPM/XPM y
        \texttt{win\_pct}; scatter plots KDA vs.\ win rate por cuartiles.
  \item \textbf{An\'alisis de objetivos:} comparaci\'on de promedios de
        \texttt{roshan\_kills\_avg} y \texttt{tower\_kills\_avg} entre victorias y
        derrotas.
  \item \textbf{An\'alisis de visi\'on:} comparaci\'on de \texttt{observer\_uses\_avg}
        y \texttt{observer\_kills\_avg} entre equipos ganadores y perdedores.
  \item \textbf{Visualizaci\'on:} once gr\'aficos exploratorios que cubren anomal\'ias,
        eficiencia econ\'omica, win rates, KDA, visi\'on, objetivos, duraci\'on,
        evoluci\'on temporal y correlaciones.
\end{enumerate}

%% ============================================================
\section{Fuente de Datos}
%% ============================================================

\subsection{Fuente}

El conjunto de datos fue publicado en \textbf{Kaggle} por el usuario Teo Calvo
bajo el nombre \textit{Dota 2 Pro Players Matches 2019--2021}~\cite{kaggle2021}.
Los datos fueron extra\'idos de la \textbf{API p\'ublica de OpenDota}, que agrega
estad\'isticas de Dota~2 directamente desde los servidores de Valve. La
recolecci\'on abarca el periodo \textbf{2019-01-01 -- 2021-06-18} e incluye
exclusivamente torneos y ligas profesionales. El procesamiento y estructuraci\'on
del dataset siguieron la metodolog\'ia descrita en el paper de referencia
\cite{paper2021}, que propone un marco para el an\'alisis de rendimiento de equipos
profesionales de Dota~2 a partir de ventanas hist\'oricas de estad\'isticas.

El conocimiento involucrado proviene del \'area de \textit{deportes electr\'onicos
(esports)} y miner\'ia de datos aplicada a videojuegos competitivos. El problema
computacional central es doble: \textbf{clasificaci\'on binaria} (\textbf{¿tiene
Radiant ventaja antes de jugar?}) y \textbf{recuperaci\'on de informaci\'on}
(b\'usqueda de partidas similares por KNN). Cada variable fue capturada porque
representa una dimensi\'on real de rendimiento hist\'orico --- econ\'omica, de
combate, de objetivos o de visi\'on --- con impacto demostrado en el resultado
final de las partidas.

\subsection{Descripci\'on}

\subsubsection*{a) A nivel de atributos}

\textbf{Archivo:} \texttt{tb\_pro\_players\_matches.csv} \\
\textbf{Tama\~no:} $\sim$50~MB \quad \textbf{Filas:} 47\,150 \quad
\textbf{Columnas:} 317 \quad \textbf{Periodo:} 2019-01-01 -- 2021-06-18

Las 317 columnas se organizan en tres grupos funcionales:

\begin{center}
\begin{tabular}{lrp{7cm}}
  \toprule
  \textbf{Grupo} & \textbf{N.º columnas} & \textbf{Descripci\'on} \\
  \midrule
  Identificadores          & 3   & \texttt{match\_id}, \texttt{dt\_match}, \texttt{radiant\_win} \\
  Estad\'isticas hist\'oricas & 72  & 36 m\'etricas $\times$ 2 equipos (\texttt{\_r} / \texttt{\_d}) \\
  Pick rate de h\'eroes    & 242 & 121 h\'eroes $\times$ 2 equipos (\texttt{\_r} / \texttt{\_d}) \\
  \bottomrule
\end{tabular}
\end{center}

\textbf{Tipos de datos:}
\begin{itemize}
  \item \textbf{Entero} (\texttt{int64}): \texttt{match\_id} --- ID \'unico de la
        partida en servidores Valve.
  \item \textbf{Objeto} (\texttt{object}): \texttt{dt\_match} --- fecha y hora exacta
        (UTC).
  \item \textbf{Booleano}: \texttt{radiant\_win} --- \texttt{True} = gan\'o Radiant,
        \texttt{False} = gan\'o Dire.
  \item \textbf{Flotante} (\texttt{float64}): todas las m\'etricas hist\'oricas
        (promedios sobre la ventana) y los pick rates de h\'eroes.
\end{itemize}

\textbf{Cuadro resumen de atributos principales:}

\begin{center}
\renewcommand{\arraystretch}{1.25}
\begin{longtable}{p{4.2cm}p{3.8cm}p{1.4cm}p{1.2cm}p{2.5cm}}
  \toprule
  \textbf{Atributo} & \textbf{Descripci\'on} & \textbf{Tipo} & \textbf{Nulos} & \textbf{Rango t\'ipico} \\
  \midrule
  \endhead
  \texttt{match\_id}
    & ID \'unico de la partida (Valve)
    & int64 & 0 & --- \\
  \texttt{dt\_match}
    & Fecha/hora exacta (UTC)
    & object & 0 & 2019--2021 \\
  \texttt{radiant\_win}
    & Resultado (True/False)
    & bool & 0 & \{T, F\} \\
  \texttt{freq\_r/d}
    & N.º partidas en la ventana hist\'orica
    & float64 & $\sim$6\,\% & 1 -- 323 \\
  \texttt{recencia\_r/d}
    & D\'ias desde \'ultima partida del equipo
    & float64 & $\sim$6\,\% & 0 -- 184 \\
  \texttt{win\_pct\_r/d}
    & Tasa de victorias hist\'oricas (0--1)
    & float64 & $\sim$6\,\% & 0.35 -- 0.75 \\
  \texttt{kda\_avg\_r/d}
    & Ratio (Kills + Asist) / Muertes
    & float64 & $\sim$6\,\% & 1.5 -- 6.0 \\
  \texttt{gold\_per\_min\_avg\_r/d}
    & Oro por minuto (GPM)
    & float64 & $\sim$6\,\% & 350 -- 600 \\
  \texttt{xp\_per\_min\_avg\_r/d}
    & Experiencia por minuto (XPM)
    & float64 & $\sim$6\,\% & 350 -- 700 \\
  \texttt{hero\_damage\_avg\_r/d}
    & Da\~no total a h\'eroes enemigos
    & float64 & $\sim$6\,\% & 8\,000 -- 30\,000 \\
  \texttt{tower\_kills\_avg\_r/d}
    & Torres enemigas destruidas (prom.)
    & float64 & $\sim$6\,\% & 2 -- 8 \\
  \texttt{tower\_damage\_avg\_r/d}
    & Da\~no infligido a torres
    & float64 & $\sim$6\,\% & 5\,000 -- 25\,000 \\
  \texttt{roshan\_kills\_avg\_r/d}
    & Kills de Roshan por partida
    & float64 & $\sim$6\,\% & 0.0 -- 1.5 \\
  \texttt{hero\_kills\_avg\_r/d}
    & H\'eroes enemigos eliminados
    & float64 & $\sim$6\,\% & 10 -- 35 \\
  \texttt{deaths\_avg\_r/d}
    & Muertes propias del equipo
    & float64 & $\sim$6\,\% & 8 -- 30 \\
  \texttt{assists\_avg\_r/d}
    & Asistencias en kills
    & float64 & $\sim$6\,\% & 20 -- 60 \\
  \texttt{kills\_per\_min\_avg\_r/d}
    & Ritmo de generaci\'on de kills
    & float64 & $\sim$6\,\% & 0.3 -- 0.8 \\
  \texttt{last\_hits\_avg\_r/d}
    & Last hits (farm de carril)
    & float64 & $\sim$6\,\% & 300 -- 900 \\
  \texttt{denies\_avg\_r/d}
    & Denies (mec\'anica exclusiva de Dota~2)
    & float64 & $\sim$6\,\% & 30 -- 120 \\
  \texttt{lane\_kills\_avg\_r/d}
    & Kills en fase de carril ($\sim$10 min)
    & float64 & $\sim$6\,\% & 2 -- 10 \\
  \texttt{firstblood\_claimed\_avg\_r/d}
    & Proporci\'on de partidas con First Blood
    & float64 & $\sim$6\,\% & 0.35 -- 0.65 \\
  \texttt{observer\_uses\_avg\_r/d}
    & Wards observer propias colocadas
    & float64 & $\sim$6\,\% & 5 -- 20 \\
  \texttt{observer\_kills\_avg\_r/d}
    & Wards observer enemigas destruidas
    & float64 & $\sim$6\,\% & 3 -- 15 \\
  \texttt{sentry\_uses\_avg\_r/d}
    & Wards sentry propias colocadas
    & float64 & $\sim$6\,\% & 4 -- 20 \\
  \texttt{sentry\_kills\_avg\_r/d}
    & Wards sentry enemigas destruidas
    & float64 & $\sim$6\,\% & 3 -- 15 \\
  \texttt{buyback\_count\_avg\_r/d}
    & Compras de revivir por partida
    & float64 & $\sim$6\,\% & 0.3 -- 2.0 \\
  \texttt{level\_avg\_r/d}
    & Nivel promedio del equipo al final
    & float64 & $\sim$6\,\% & 18 -- 28 \\
  \texttt{actions\_per\_min\_avg\_r/d}
    & Acciones por minuto (APM)
    & float64 & $\sim$6\,\% & 150 -- 350 \\
  \texttt{neutral\_kills\_avg\_r/d}
    & Kills de neutrales en jungla
    & float64 & $\sim$6\,\% & 50 -- 300 \\
  \texttt{duration\_avg\_win\_r/d}
    & Duraci\'on promedio en victorias (min)
    & float64 & $\sim$6\,\% & 30 -- 50 \\
  \texttt{duration\_avg\_lose\_r/d}
    & Duraci\'on promedio en derrotas (min)
    & float64 & $\sim$6\,\% & 32 -- 55 \\
  \texttt{hero\_<ID>\_avg\_r/d}
    & Pick rate del h\'eroe (proporci\'on)
    & float64 & $\sim$6\,\% & 0.0 -- 1.0 \\
  \bottomrule
\end{longtable}
\end{center}

\noindent\textbf{Clasificaci\'on de variables:}
\begin{itemize}
  \item \texttt{radiant\_win}: \textbf{categ\'orica nominal binaria} (variable
        objetivo).
  \item M\'etricas de rendimiento (\texttt{kda\_avg}, \texttt{gold\_per\_min\_avg},
        etc.): \textbf{cuantitativas continuas}.
  \item \texttt{hero\_<ID>\_avg}: \textbf{cuantitativa continua}, interpretable
        como proporci\'on (rango $[0,1]$).
  \item \texttt{freq\_r/d}, \texttt{recencia\_r/d}: \textbf{cuantitativas
        discretas}, usadas como metadatos de confiabilidad.
\end{itemize}

\noindent\textbf{Valores nulos:} el $\sim$\textbf{6\,\%} del dataset (2\,819
partidas) presenta \texttt{NaN} en todas las columnas de estad\'isticas.
Corresponden a partidas no completadas por abandono de jugadores, partidas de
torneo jugadas simult\'aneamente que comparten resultado y rendimiento, o equipos
sin historial en la API. Se gestionan mediante imputaci\'on con la media columnar.

\noindent\textbf{Medidas de dispersi\'on --- ventana hist\'orica:}

\begin{center}
\begin{tabular}{lrrrr}
  \toprule
  \textbf{Variable} & \textbf{Media} & \textbf{Mediana} & \textbf{M\'ax.} & \textbf{Desv. Est.} \\
  \midrule
  \texttt{freq\_r} (partidas)   & 74.9 & 63 & 323 & $\sim$52 \\
  \texttt{freq\_d} (partidas)   & 74.7 & 62 & 323 & $\sim$52 \\
  \texttt{recencia\_r} (d\'ias) &  5.2 &  2 & 184 & $\sim$12 \\
  \texttt{recencia\_d} (d\'ias) &  5.1 &  2 & 184 & $\sim$12 \\
  \bottomrule
\end{tabular}
\end{center}

\subsubsection*{b) A nivel de registros}

Cada registro representa \textbf{una partida profesional de Dota~2} identificada por
un \texttt{match\_id} \'unico --- sin duplicados. Los registros est\'an
\textbf{etiquetados} mediante \texttt{radiant\_win}: \texttt{True} indica victoria
de Radiant y \texttt{False} victoria de Dire.

No hay granularidad geogr\'afica. La granularidad temporal es \textbf{a nivel de
partida individual} con fecha y hora exacta en UTC. La distribuci\'on temporal es:

\begin{center}
\begin{tabular}{lrr}
  \toprule
  \textbf{A\~no} & \textbf{Partidas} & \textbf{\% del total} \\
  \midrule
  2019              & 17\,578 & 37.3\,\% \\
  2020              & 22\,376 & 47.4\,\% \\
  2021 (hasta jun.) &  7\,196 & 15.3\,\% \\
  \midrule
  \textbf{Total}   & \textbf{47\,150} & \textbf{100\,\%} \\
  \bottomrule
\end{tabular}
\end{center}

Resultados globales: \textbf{Radiant gana 52.1\,\%} y Dire gana 47.9\,\%.
Esta diferencia es un fen\'omeno conocido en Dota~2 relacionado con la geometr\'ia
del mapa y la posici\'on de los objetivos secundarios.

\subsubsection*{c) Relaci\'on entre atributos}

Existe una relaci\'on directa entre \texttt{win\_pct\_r/d} y las m\'etricas de
rendimiento: los equipos con mayor tasa de victorias hist\'orica tienden a registrar
KDA, GPM y XPM m\'as elevados, dado que las partidas ganadas acumulan m\'as kills,
asistencias y eficiencia econ\'omica. La relaci\'on \texttt{tower\_kills\_avg} --
\texttt{win\_pct} es especialmente relevante porque destruir torres es el objetivo
primario del juego. Las variables \texttt{freq\_r/d} y \texttt{recencia\_r/d} son
metadatos de confiabilidad y no miden rendimiento directamente.

\textbf{Terminolog\'ia especial del dataset:}
\begin{itemize}
  \item \textbf{GPM} (\textit{Gold Per Minute}): ritmo de generaci\'on de oro;
        indicador primario de eficiencia econ\'omica.
  \item \textbf{XPM} (\textit{Experience Per Minute}): ritmo de ganancia de
        experiencia; determina el nivel del h\'eroe.
  \item \textbf{KDA}: ratio (Kills + Asistencias) / Muertes de equipo.
        KDA $>$ 4 es excelente; KDA $<$ 1.5 suele correlacionar con derrotas.
  \item \textbf{Roshan}: boss neutral de alto valor; matarlo otorga el Aegis
        (vida extra) y gran recompensa de oro/XP.
  \item \textbf{Ventana hist\'orica}: las $N$ partidas previas usadas para
        calcular los promedios; \texttt{freq\_r/d} indica su tama\~no.
  \item \textbf{Pick rate}: fracci\'on de partidas en que un h\'eroe fue elegido
        dentro de la ventana hist\'orica.
  \item \textbf{Last hits / Denies}: last hit = \'ultimo golpe a un creep para
        obtener oro; deny = negar ese oro al enemigo (mec\'anica exclusiva de
        Dota~2).
  \item \textbf{TP Scroll}: pergamino de teletransporte a torre aliada; permite
        rotaciones globales.
  \item \textbf{First Blood}: primera kill de la partida; otorga $\sim$200 de
        oro extra al equipo que la consigue.
\end{itemize}

\subsection{Formato}

El dataset se encontr\'o en un \'unico archivo de texto plano con separaci\'on por
comas (\textbf{CSV}): \texttt{tb\_pro\_players\_matches.csv} ($\sim$50~MB).
Contiene 47\,150 filas y 317 columnas. Los tipos de datos originales son:
entero (\texttt{match\_id}), booleano (\texttt{radiant\_win}), flotante (m\'etricas
y pick rates) y texto (\texttt{dt\_match}).

\subsection{Transformaciones}

Las transformaciones aplicadas para convertir los datos en un formato utilizable
para el modelo KNN fueron:

\begin{table}[h!]
\centering
\renewcommand{\arraystretch}{1.3}
\begin{tabular}{@{}c p{4.5cm} p{7cm}@{}}
  \toprule
  \textbf{Paso} & \textbf{T\'ecnica} & \textbf{Justificaci\'on} \\
  \midrule
  1 & C\'omputo de flags \texttt{\_has\_heroes}
    & Detectar autom\'aticamente las 2\,819 filas sin historial de h\'eroes
      (abandonos, partidas simult\'aneas o equipos sin registro) \\
  2 & Selecci\'on de 66 columnas de estad\'isticas
    & Se excluyen \texttt{recencia}, \texttt{freq}, \texttt{lane\_efficiency\_avg}
      y los 242 pick rates: la alta dimensionalidad degrada la distancia euclidiana \\
  3 & Imputaci\'on de NaN con la media columnar
    & El $\sim$6\,\% de partidas sin historial no puede descartarse;
      la media es la mejor estimaci\'on disponible \\
  4 & Construcci\'on de \'indice \textbf{BallTree}
    & B\'usqueda eficiente de vecinos con m\'etrica euclidiana sobre 47\,150 registros \\
  5 & Consulta KNN con el mismo scaler
    & Garantiza comparabilidad entre el vector de consulta y el \'indice \\
  \bottomrule
\end{tabular}
\end{table}

%% ============================================================
\section{Exploración}
%% ============================================================
 
Se realizaron \textbf{once gráficos exploratorios} para investigar las hipótesis
planteadas y describir el dataset desde múltiples dimensiones: calidad de datos,
balance de resultados, eficiencia económica, meta de héroes, rendimiento de
combate, control de visión, objetivos, duración, evolución temporal y
correlaciones con la variable objetivo. A continuación se presenta cada
visualización con su descripción detallada y los hallazgos obtenidos.
 
% ------------------------------------------------------------------
\subsection{Gráfico 1 --- Distribución de partidas válidas e inválidas por año}
 
\includegraphics[width=0.95\linewidth]{grafico1.png}
 
\textbf{Descripción técnica:} Gráfico de barras apiladas que desglosa, para
cada año del periodo analizado (2019, 2020 y 2021), la cantidad de partidas
\textbf{válidas} (aquellas que tienen al menos un héroe con pick rate $>$ 0 en
ambos equipos, \texttt{\_has\_heroes = True}) frente a las partidas
\textbf{anómalas} (\texttt{\_has\_heroes = False}). Las barras azules representan
las partidas válidas y las rojas las anomalías. El porcentaje de anomalía
por año se etiqueta en rojo sobre la barra correspondiente.
 
\textbf{Metodología:} Se calcula la columna \texttt{\_has\_heroes} como la
intersección lógica de \texttt{\_has\_heroes\_r} y \texttt{\_has\_heroes\_d},
donde cada flag indica si al menos un héroe del equipo tiene valor $>$ 0 en la
ventana histórica. Las partidas anómalas corresponden a equipos sin historial
registrado en OpenDota (debutantes, datos faltantes en la API).
 
\textbf{Hallazgos:}
\begin{itemize}
  \item El total de anomalías es de $\sim$2\,819 partidas, equivalente a
        aproximadamente el \textbf{6\,\%} del dataset.
  \item La tasa de anomalía es relativamente estable entre años, lo que
        indica que el problema de datos faltantes no se concentra en ningún
        periodo específico.
  \item El año 2020 concentra el mayor volumen absoluto de partidas (47.4\,\%
        del total), pero no la mayor proporción de anomalías.
  \item Estas anomalías se excluyen del análisis de héroes y del modelo KNN,
        pero se conservan para el análisis de estadísticas globales con
        imputación de la media.
\end{itemize}
 
% ------------------------------------------------------------------
\subsection{Gráfico 2 --- Win Rate global: Radiant vs.\ Dire}
 
\includegraphics[width=0.55\linewidth]{grafico2.png}
 
\textbf{Descripción técnica:} Gráfico tipo \textit{donut} (anillo) que muestra
la distribución porcentual de victorias entre Radiant y Dire sobre las 44\,331
partidas válidas (con historial de héroes). Los porcentajes se calculan directamente sobre la columna booleana
\texttt{radiant\_win}.
 
\textbf{Hallazgos:}
\begin{itemize}
  \item \textbf{Radiant gana el 52.1\,\%} de las partidas válidas y Dire
        el 47.9\,\%.
  \item Esta diferencia de $\sim$4.2 puntos porcentuales es estadísticamente
        significativa sobre una muestra de más de 44\,000 partidas.
  \item El sesgo de mapa es un fenómeno bien documentado en Dota~2 competitivo:
        la base de Radiant se ubica en la esquina inferior izquierda, con acceso
        más directo al objetivo de Roshan y a los creeps de jungla de mayor valor.
  \item Este desequilibrio implica que cualquier modelo predictivo debe considerar
        el equipo (Radiant vs.\ Dire) como variable relevante.
\end{itemize}
 
% ------------------------------------------------------------------
\subsection{Gráfico 3 --- GPM y XPM: Radiant vs.\ Dire}
 
\includegraphics[width=0.85\linewidth]{grafico3.png}
 
\textbf{Descripción técnica:} Barras agrupadas que comparan el promedio global
de \texttt{gold\_per\_min\_avg} (GPM) y \texttt{xp\_per\_min\_avg} (XPM) entre
Radiant y Dire calculados sobre todas las partidas válidas. Cada par de barras
representa una métrica. El valor promedio se etiqueta sobre cada barra.
 
\textbf{Metodología:} Se calculan los promedios directamente de las columnas
\texttt{gold\_per\_min\_avg\_r}, \texttt{gold\_per\_min\_avg\_d},
\texttt{xp\_per\_min\_avg\_r} y \texttt{xp\_per\_min\_avg\_d} del subconjunto
\texttt{df\_valid}.
 
\textbf{Hallazgos:}
\begin{itemize}
  \item Radiant presenta un GPM ligeramente superior al de Dire, consistente con
        su mayor tasa de victorias (los equipos ganadores acumulan más oro).
  \item El XPM sigue el mismo patrón: mayor para Radiant en promedio.
  \item La diferencia en GPM y XPM es pequeña en términos absolutos cuando
        se promedia sobre todo el dataset, porque el resultado de la partida
        no está controlado --- ambos equipos aparecen como ganadores y
        perdedores en diferentes filas.
  \item Esto refuerza la hipótesis 1: la ventaja económica de Radiant no es
        solo estructural de mapa, sino que se refleja en las métricas
        históricas de eficiencia.
\end{itemize}
 
% ------------------------------------------------------------------
\subsection{Gráfico 4 --- Win Rate de los Top 50 Héroes más elegidos}
 
\includegraphics[width=0.95\linewidth]{grafico4.png}
 
\textbf{Descripción técnica:} Gráfico de 50 barras verticales que muestra el
win rate individual de cada uno de los 50 héroes más utilizados en el dataset.
La \textbf{frecuencia} de un héroe se calcula como la \textbf{suma acumulada de
sus pick rates} (no como conteo de partidas), lo que pondera habilidades elegidas
consistentemente sobre muchas partidas. Las barras se ordenan de mayor a menor win
rate para una lectura fluida, y se colorean según tres umbrales: verde
($\geq$55\,\%), rojo ($\leq$45\,\%) y dorado (rango intermedio). Líneas de
referencia horizontales a 45\,\%, 50\,\% y 55\,\%.
 
\textbf{Metodología:}
\begin{enumerate}
  \item Para cada héroe $h$, se suma \texttt{hero\_h\_avg\_r.fillna(0)} sobre
        todas las filas ($\text{freq\_r}$) y equivalente para Dire.
  \item Las victorias se cuentan: si Radiant tiene el héroe y ganó
        (\texttt{radiant\_win = True}) suma como victoria; si está en Dire y
        perdió Radiant (\texttt{radiant\_win = False}), también.
  \item $\text{WR}(h) = \text{wins}(h) / \text{partidas}(h) \times 100$.
  \item Solo se incluyen héroes con $\geq$50 partidas de aparición.
\end{enumerate}
 
\textbf{Hallazgos:}
\begin{itemize}
  \item Existe un rango relativamente estrecho de win rates (aproximadamente
        43--57\,\%) en el ámbito profesional, lo que indica que el balance del
        meta es alto.
  \item Los héroes con win rate $\geq$55\,\% son los considerados ``meta'' o
        sobrevaluados en el parche analizado.
  \item Los héroes con win rate $\leq$45\,\% son frecuentes a pesar de su baja
        efectividad, lo que sugiere que su uso responde a estrategias específicas
        o al drafting del equipo contrario.
  \item La metodología de suma de pick rates (en lugar de conteo de filas)
        garantiza que el ranking sea idéntico al de la aplicación web, donde
        un héroe elegido con alta tasa en pocas partidas pesa más que uno
        elegido puntualmente en muchas.
\end{itemize}
 
% ------------------------------------------------------------------
\subsection{Gráfico 5 --- KDA: ganadores vs.\ perdedores}
 
\includegraphics[width=0.85\linewidth]{grafico5.png}
 
\textbf{Descripción técnica:} Histograma de densidad superpuesto que compara
la distribución del KDA histórico (\texttt{kda\_avg}) entre los equipos que
ganaron y los que perdieron en cada partida. Para cada partida válida, el equipo
ganador es Radiant cuando \texttt{radiant\_win = True}, o Dire en caso contrario.
Los valores se recortan al rango $[0.5,\, 8]$ para eliminar outliers extremos.
Las líneas verticales punteadas indican la media de cada distribución.
 
\textbf{Metodología:} Se extraen los KDA de ganadores concatenando
\texttt{kda\_avg\_r} (cuando Radiant gana) con \texttt{kda\_avg\_d} (cuando Dire
gana), y análogamente para perdedores. Se usa la función de densidad (no
frecuencia absoluta) para hacer comparables distribuciones de tamaños similares.
 
\textbf{Hallazgos:}
\begin{itemize}
  \item Los equipos \textbf{ganadores} presentan una distribución de KDA
        desplazada hacia la derecha respecto a los perdedores, con media
        significativamente mayor.
  \item La superposición de distribuciones indica que el KDA no es un predictor
        perfecto: un equipo puede ganar con KDA bajo (por objetivos, no por
        combate) o perder con KDA alto (si cede torres).
  \item Apoya la hipótesis 2: mayor KDA histórico se asocia con mayor
        probabilidad de victoria, aunque no de forma determinista.
  \item La cola derecha de los ganadores es más larga, indicando que los equipos
        con rendimiento excepcional de combate casi siempre ganan.
\end{itemize}
 
% ------------------------------------------------------------------
\subsection{Gráfico 6 --- Control de visión: wards colocadas y destruidas}
 
\includegraphics[width=0.85\linewidth]{grafico6.png}
 
\textbf{Descripción técnica:} Barras agrupadas que comparan cuatro métricas
de control de visión entre equipos ganadores y perdedores: wards
\textbf{observer} colocadas (\texttt{observer\_uses\_avg}), wards observer
\textbf{destruidas} (\texttt{observer\_kills\_avg}), wards \textbf{sentry}
colocadas (\texttt{sentry\_uses\_avg}) y wards sentry \textbf{destruidas}
(\texttt{sentry\_kills\_avg}). Cada par de barras agrupa ganador (verde) vs.\
perdedor (rojo) por métrica.
 
\textbf{Metodología:} Se utiliza la función \texttt{split\_win\_lose(col\_r,
col\_d)} que separa los valores por resultado, concatenando el equipo correcto
(Radiant o Dire) según \texttt{radiant\_win}. Los promedios se calculan sobre
todas las partidas válidas.
 
\textbf{Hallazgos:}
\begin{itemize}
  \item Los equipos ganadores colocan \textbf{más} wards observer y sentry
        que los perdedores, indicando mayor inversión en información.
  \item Los ganadores también \textbf{destruyen más} wards enemigas,
        lo que se conoce como ``desvisión'' activa.
  \item La diferencia es más pronunciada en observer kills que en sentry kills,
        sugiriendo que negar la visión pasiva del enemigo es más diferenciador.
  \item Apoya la hipótesis 3: la guerra de visión es un diferenciador
        estadísticamente real en el circuito profesional.
  \item Equipos con mejor visión toman mejores decisiones de mapa (ganks,
        rotaciones, objetivos), lo que se traduce directamente en ventaja de
        recursos.
\end{itemize}
 
% ------------------------------------------------------------------
\subsection{Gráfico 7 --- Torres vs.\ Kills: ¿qué predice mejor la victoria?}
 
\includegraphics[width=0.85\linewidth]{grafico7.png}
 
\textbf{Descripción técnica:} Gráfico de barras horizontales que cuantifica
la \textbf{diferencia relativa} entre ganadores y perdedores para cinco tipos de
objetivos: torres destruidas (\texttt{tower\_kills\_avg}), kills de héroes
(\texttt{hero\_kills\_avg}), daño a torres (\texttt{tower\_damage\_avg}), kills
de Roshan (\texttt{roshan\_kills\_avg}) y kills de neutrales
(\texttt{neutral\_kills\_avg}). La diferencia relativa se calcula como:
 
$$\Delta_{\%}(m) = \frac{\bar{m}_\text{ganador} - \bar{m}_\text{perdedor}}{\bar{m}_\text{perdedor}} \times 100$$
 
Las barras verdes indican que el ganador supera al perdedor. El valor se etiqueta
al extremo de cada barra.
 
\textbf{Hallazgos:}
\begin{itemize}
  \item \textbf{Las torres destruidas} presentan la mayor diferencia relativa
        entre ganadores y perdedores, confirmando la hipótesis 5: las torres
        son el predictor más fuerte del resultado.
  \item El \textbf{daño a torres} es el segundo diferenciador más grande,
        lo que indica que el trabajo de siegeo (asedio a estructuras) es clave.
  \item Las \textbf{kills de héroes} tienen una diferencia positiva pero menor,
        lo que sugiere que ganar combates no garantiza la victoria si no se
        traduce en objetivos estructurales.
  \item Los \textbf{kills de Roshan} muestran cierta diferencia, pero son
        relativamente bajos en valor absoluto (Roshan no siempre se mata).
  \item Los \textbf{kills de neutrales} son los menos diferenciadores, pues
        ambos equipos farman la jungla de manera similar independientemente del
        resultado.
\end{itemize}
 
% ------------------------------------------------------------------
\subsection{Gráfico 8 --- Duración de partidas: ganando vs.\ perdiendo}
 
\includegraphics[width=0.95\linewidth]{grafico8.png}
 
\textbf{Descripción técnica:} Figura de dos paneles. El panel izquierdo es un
histograma de densidad superpuesto que compara la distribución de
\texttt{duration\_avg\_win\_r} (duración promedio histórica en victorias de
Radiant) contra \texttt{duration\_avg\_lose\_r} (en derrotas). El panel derecho
presenta un boxplot comparativo entre ambas distribuciones. Los valores están
en minutos (la columna original está en segundos y se divide entre 60).
 
\textbf{Hallazgos:}
\begin{itemize}
  \item Los equipos de Radiant tienen una duración promedio de victoria
        \textbf{menor} que su duración promedio de derrota.
  \item Esto confirma un principio conocido en Dota~2: los equipos \textbf{ganan
        más rápido} de lo que pierden, porque el equipo ganador ejecuta la
        victoria cuando tiene ventaja, mientras que el equipo perdedor alarga el
        juego buscando una remontada (\textit{comeback}).
  \item El boxplot muestra mayor varianza en las derrotas, lo que indica que la
        duración de los juegos perdidos es más impredecible.
  \item Partidas muy largas ($>$50 min) suelen corresponder a derrotas o a
        partidas muy parejas donde la ventaja se escapó.
  \item Relevante para la hipótesis 4: en partidas largas, Roshan se mata
        más veces y su impacto relativo aumenta.
\end{itemize}
 
% ------------------------------------------------------------------
\subsection{Gráfico 9 --- GPM por año: ganadores vs.\ perdedores}
 
\includegraphics[width=0.85\linewidth]{grafico9.png}
 
\textbf{Descripción técnica:} Barras agrupadas que muestran, para cada año
del periodo analizado (2019, 2020, 2021), el GPM (\texttt{gold\_per\_min\_avg})
promedio de los equipos ganadores frente a los perdedores. El análisis se
realiza sobre el subconjunto de partidas válidas de cada año. Cada par de
barras representa un año, con barra verde para ganadores y roja para perdedores.
 
\textbf{Metodología:} Se filtra \texttt{df\_valid} por \texttt{\_year} y se
separan los GPM de ganadores y perdedores usando \texttt{split\_win\_lose} por
año, concatenando Radiant (cuando \texttt{radiant\_win = True}) y Dire (cuando
\texttt{radiant\_win = False}).
 
\textbf{Hallazgos:}
\begin{itemize}
  \item La brecha de GPM entre ganadores y perdedores se mantiene
        \textbf{consistente en los tres años}, lo que confirma que la eficiencia
        económica es un predictor estable del resultado a lo largo del tiempo.
  \item El GPM promedio general aumentó entre 2019 y 2021, reflejando la
        evolución del meta profesional hacia juegos con mayor ritmo de farm.
  \item La diferencia absoluta entre ganadores y perdedores es de
        aproximadamente 30--50 GPM en todos los años.
  \item Esto valida la hipótesis 1 de forma temporal: la eficiencia económica
        no es un artefacto de un parche específico, sino una constante
        estructural del juego profesional.
\end{itemize}
 
% ------------------------------------------------------------------
\subsection{Gráfico 10 --- Tamaño de ventana histórica y estabilidad del win rate}
 
\includegraphics[width=0.85\linewidth]{grafico10.png}
 
\textbf{Descripción técnica:} Gráfico combinado con doble eje Y. El eje
izquierdo (barras cianas) muestra el win rate promedio (\texttt{win\_pct\_r})
por tramo de tamaño de ventana histórica (\texttt{freq\_r}), agrupado en
6 buckets: 1--10, 11--30, 31--60, 61--100, 101--200 y 201+. Las barras de error
representan la desviación estándar del win rate dentro de cada bucket. El eje
derecho (línea dorada) muestra el número de partidas en cada bucket.
 
\textbf{Hallazgos:}
\begin{itemize}
  \item Los equipos con \textbf{ventana histórica pequeña} (1--10 partidas)
        presentan una \textbf{mayor desviación estándar} en su win rate,
        lo que indica que sus estadísticas son menos representativas.
  \item A medida que aumenta \texttt{freq\_r}, la desviación estándar del
        win rate se reduce: con más historial, las métricas convergen a
        valores más confiables.
  \item El bucket 201+ tiene el menor error, aunque también el menor número
        de partidas (equipos muy activos).
  \item Este hallazgo justifica el uso de la media columnar para imputar valores
        faltantes: los equipos sin historial serán asignados al promedio global,
        que es la mejor estimación disponible.
  \item Es relevante para el modelo KNN: distancias calculadas sobre equipos con
        poca historia son menos confiables que las de equipos con $>$60 partidas
        históricas.
\end{itemize}
 
% ------------------------------------------------------------------
\subsection{Gráfico 11 --- Correlación de estadísticas con la victoria}
 
\includegraphics[width=0.95\linewidth]{grafico11.png}
 
\textbf{Descripción técnica:} Gráfico de barras horizontales que muestra el
\textbf{coeficiente de correlación de Pearson} entre 17 estadísticas
seleccionadas de Radiant y la variable objetivo \texttt{radiant\_win}. Las
estadísticas incluyen: \texttt{win\_pct}, \texttt{kda\_avg},
\texttt{gold\_per\_min\_avg}, \texttt{xp\_per\_min\_avg}, \texttt{hero\_kills\_avg},
\texttt{deaths\_avg}, \texttt{assists\_avg}, \texttt{tower\_kills\_avg},
\texttt{tower\_damage\_avg}, \texttt{roshan\_kills\_avg},
\texttt{observer\_uses\_avg}, \texttt{sentry\_uses\_avg}, \texttt{last\_hits\_avg},
\texttt{denies\_avg}, \texttt{neutral\_kills\_avg},
\texttt{firstblood\_claimed\_avg} y \texttt{buyback\_count\_avg}. Las barras se
colorean en verde si la correlación es positiva y en rojo si es negativa, y se
ordenan de mayor a menor correlación.
 
\textbf{Metodología:} Se calcula la matriz de correlación de Pearson sobre
\texttt{df\_valid} para las 17 features seleccionadas más \texttt{radiant\_win},
y se extrae la columna correspondiente a \texttt{radiant\_win}.
 
\textbf{Hallazgos:}
\begin{itemize}
  \item \texttt{win\_pct\_r} tiene la \textbf{correlación más alta} con
        \texttt{radiant\_win}, lo que es esperado: la tasa de victorias histórica
        predice victorias futuras.
  \item \texttt{kda\_avg\_r}, \texttt{gold\_per\_min\_avg\_r} y
        \texttt{xp\_per\_min\_avg\_r} tienen correlaciones positivas moderadas,
        confirmando las hipótesis 1 y 2.
  \item \texttt{deaths\_avg\_r} tiene la \textbf{correlación negativa más
        fuerte}: más muertes históricas predicen menor probabilidad de ganar.
  \item \texttt{tower\_kills\_avg\_r} tiene mayor correlación que
        \texttt{hero\_kills\_avg\_r}, apoyando la hipótesis 5.
  \item Variables como \texttt{buyback\_count\_avg\_r} tienen correlaciones
        prácticamente nulas, lo que indica que el uso de buybacks no discrimina
        ganadores de perdedores.
  \item Este gráfico sirve como \textbf{mapa de importancia de variables} para
        orientar futuros modelos de clasificación supervisados sobre este dataset.
\end{itemize}

%% ============================================================
\section{Conclusi\'on}
%% ============================================================

\subsection*{Hip\'otesis 1 --- Eficiencia econ\'omica como predictor de victoria}

\textbf{Conclusi\'on:} Los gr\'aficos 3, 9 y 11 confirman que el GPM y el XPM son
predictores positivos y consistentes de la victoria en Dota~2 profesional. La
brecha de GPM entre ganadores y perdedores se mantiene estable en los tres a\~nos
analizados ($\sim$30--50 GPM de diferencia). El GPM ocupa un lugar prominente en
el ranking de correlaci\'on de Pearson con \texttt{radiant\_win}. Los equipos con
mayor eficiencia econ\'omica adquieren \'items m\'as poderosos antes que sus
rivales, generando una ventaja de combate compuesta y sostenida.

\subsection*{Hip\'otesis 2 --- KDA hist\'orico y tasa de victorias}

\textbf{Conclusi\'on:} El gr\'afico 5 muestra una separaci\'on clara de las
distribuciones de KDA entre ganadores y perdedores. Los equipos con
\texttt{kda\_avg} $>$ 3.5 presentan win rates notablemente superiores. El KDA
captura sint\'eticamente el rendimiento de combate: m\'as kills y asistencias con
menos muertes implica mayor control de los enfrentamientos y mayor acumulaci\'on
de recursos para el equipo. El gr\'afico 11 confirma que \texttt{kda\_avg\_r}
es la segunda variable con mayor correlaci\'on positiva tras \texttt{win\_pct\_r}.

\subsection*{Hip\'otesis 3 --- Control de visi\'on}

\textbf{Conclusi\'on:} El gr\'afico 6 muestra diferencias sistem\'aticas en todas
las m\'etricas de visi\'on entre ganadores y perdedores. Los equipos ganadores
invierten m\'as en wards (observer y sentry) y destruyen m\'as wards enemigas. La
``guerra de visi\'on'' es un diferenciador real en el circuito profesional: la
informaci\'on del mapa permite decisiones t\'acticas superiores (ganks, rotaciones,
toma de objetivos) que se traducen en ventajas econ\'omicas y estructurales.

\subsection*{Hip\'otesis 4 --- Roshan y duraci\'on}

\textbf{Conclusi\'on:} El gr\'afico 7 muestra que los kills de Roshan tienen una
diferencia relativa positiva entre ganadores y perdedores, aunque menor que la de
torres. El gr\'afico 8 confirma que en partidas largas (donde Roshan se mata m\'as
veces) la varianza del resultado aumenta. Matar a Roshan otorga el Aegis (vida
extra) y una importante recompensa de oro y XP, lo que permite ejecutar combates
de alto riesgo en las etapas tard\'ias de la partida.

\subsection*{Hip\'otesis 5 --- Torres vs.\ kills de h\'eroes}

\textbf{Conclusi\'on:} El gr\'afico 7 demuestra que \texttt{tower\_kills\_avg}
tiene la mayor diferencia relativa entre ganadores y perdedores ($>$50\,\%),
superando a \texttt{hero\_kills\_avg} ($\sim$20--30\,\%). El gr\'afico 11 confirma
este resultado con correlaciones de Pearson: las torres son m\'as predictivas del
resultado que los kills. Esto es consistente con el dise\~no del juego: el objetivo
final es destruir el Ancient, que solo es accesible tras eliminar todas las torres
defensivas.

\subsection*{Hallazgos adicionales}

\begin{itemize}
  \item \textbf{Sesgo de mapa (Gr\'afico 2):} Radiant mantiene una ventaja global
        del 52.1\,\% de manera constante, lo que debe considerarse como variable
        de control en cualquier modelo predictivo.
  \item \textbf{Calidad de datos (Gr\'afico 1):} El 6\,\% de anomal\'ias es
        manejable y se distribuye uniformemente entre a\~nos, sin comprometer la
        validez del an\'alisis.
  \item \textbf{Duraci\'on (Gr\'afico 8):} Los equipos ganan m\'as r\'apido de
        lo que pierden, lo que confirma que el equipo con ventaja ejecuta el cierre
        del juego activamente.
  \item \textbf{Estabilidad temporal (Gr\'afico 9):} Los patrones de eficiencia
        econ\'omica se mantienen consistentes en 2019--2021, lo que valida
        la generalizabilidad del modelo.
  \item \textbf{Confiabilidad del historial (Gr\'afico 10):} Las ventanas
        hist\'oricas $>$60 partidas producen estad\'isticas m\'as estables,
        lo que orienta el umbral m\'inimo para confiar en las predicciones del KNN.
  \item \textbf{Mapa de importancia (Gr\'afico 11):} El ranking de correlaciones
        orienta la selecci\'on de features para modelos supervisados futuros
        (regresi\'on log\'istica, Random Forest, XGBoost).
\end{itemize}

\subsection*{Desaf\'ios encontrados y sorpresas}

\textbf{Desaf\'ios:}
\begin{itemize}
  \item \textbf{Datos faltantes por dise\~no:} el 6\,\% de las partidas no ten\'ia
        hist\'orico de h\'eroes registrado. Esto no fue un error del CSV, sino
        una consecuencia real de la API: equipos debutantes, partidas jugadas
        simult\'aneamente que comparten estad\'isticas, y abandonos tempranos.
        El reto fue distinguir entre ``dato ausente por error'' y ``dato ausente
        porque la informaci\'on nunca existi\'o'', y decidir que la imputaci\'on
        con la media era preferible a descartar esas filas por completo.
  \item \textbf{Consistencia entre herramientas:} se detect\'o una discrepancia
        entre el ranking de h\'eroes calculado en el cuaderno de an\'alisis y el
        que mostraba la aplicaci\'on web. La causa fue sutil: el cuaderno contaba
        \textit{filas} donde aparec\'ia un h\'eroe, mientras que el \textit{backend}
        sumaba los \textit{valores num\'ericos} del pick rate. Unificar ambos
        c\'alculos requiri\'o revisar la l\'ogica completa del \textit{pipeline}.
  \item \textbf{Alta dimensionalidad:} el dataset tiene 242 columnas de
        pick rate de h\'eroes. Incluirlas todas en el modelo KNN degrada la
        m\'etrica de distancia (el llamado ``curse of dimensionality''), por lo
        que fue necesario seleccionar solo las 66 columnas de estad\'isticas
        de rendimiento m\'as relevantes.
\end{itemize}

\textbf{Sorpresas encontradas durante el an\'alisis:}
\begin{itemize}
  \item \textbf{La ventaja de Radiant es real y estable.} Se sab\'ia de
        antemano que exist\'ia cierto sesgo de mapa, pero encontrar que ese
        desequilibrio se manten\'ia en exactamente el 52.1\,\% de forma
        constante en los tres a\~nos fue notable. En 47\,150 partidas, una
        diferencia de 4 puntos porcentuales es estad\'isticamente robusta.
  \item \textbf{Las torres predicen mejor que los kills.} La hip\'otesis de
        que destruir torres es m\'as determinante que ganar combates se
        confirm\'o con una diferencia relativa del $>$50\,\% en torres frente
        al $\sim$25\,\% en kills. En un juego donde los combates son
        espect\'aculares y visibles, resultaba contraintuitivo que el ``aburrido''
        trabajo de asedio sea el verdadero predictor del resultado.
  \item \textbf{El GPM aument\'o entre 2019 y 2021.} Al analizar el GPM por
        a\~no (Gr\'afico 9) se observ\'o que el ritmo de generaci\'on de oro
        aument\'o progresivamente. Esto refleja una tendencia documentada del
        meta profesional de Dota~2 hacia estrategias de farm acelerado.
  \item \textbf{Roshan tiene un impacto m\'as d\'ebil de lo esperado.}
        Roshan es el enemigo neutro m\'as ic\'onico del juego y su muerte es
        un evento categ\'orico en toda partida. Sin embargo, su correlaci\'on
        con la victoria fue la m\'as baja entre los objetivos, probablemente
        porque ambos equipos matan a Roshan a lo largo de la partida y el
        impacto de un kill individual se diluye en el promedio hist\'orico.
\end{itemize}

\subsection*{Conclusi\'on general}

\begin{itemize}
  \item El dataset de 47\,150 partidas profesionales contiene informaci\'on
        hist\'orica rica y multidimensional que permite identificar patrones de
        rendimiento estad\'isticamente significativos.
  \item Radiant mantiene una ventaja global del 52.1\,\%, consistente con el sesgo
        de mapa documentado en Dota~2 competitivo.
  \item Las 66 features seleccionadas para el KNN capturan dimensiones reales de
        rendimiento y permiten recuperar partidas estructuralmente similares.
  \item La normalizaci\'on con StandardScaler es imprescindible dado el rango
        heterog\'eneo de las variables (GPM $\sim$400 vs.\ Roshan $\sim$0.3).
  \item Las partidas sin historial ($\sim$6\,\%) son gestionables mediante
        imputaci\'on con la media columnar sin comprometer el modelo KNN.
  \item La combinaci\'on de ventanas hist\'oricas y KNN es un enfoque s\'olido para
        el an\'alisis de esports competitivos, extensible a modelos predictivos
        de mayor complejidad.
\end{itemize}

\newpage
\section{Glosario de T\'erminos}

\textit{Esta secci\'on explica los t\'erminos del videojuego y del an\'alisis
de datos utilizados en el informe, pensada para el lector que no est\'a
familiarizado con Dota~2 ni con la miner\'ia de datos.}

\subsection*{T\'erminos de Dota 2}

\begin{description}[leftmargin=3.8cm, style=nextline, itemsep=5pt]

  \item[Dota 2 (MOBA)] Videojuego de estrategia en tiempo real de Valve
    Corporation. ``MOBA'' significa \textit{Multiplayer Online Battle Arena}:
    dos equipos de 5 jugadores se enfrentan en un mapa con el \'unico objetivo
    de destruir la base central del rival.

  \item[Radiant y Dire] Los nombres de los dos bandos en cada partida.
    Radiant ocupa la esquina inferior izquierda del mapa y Dire la superior
    derecha. En cada partida profesional se sortea qu\'e equipo juega cada bando.

  \item[Ancient] La estructura central de la base de cada equipo.
    Destruir el Ancient rival es la condici\'on de victoria. Es equivalente
    al ``Nexo'' en League of Legends o al ``cuartel general'' en los juegos
    de estrategia cl\'asicos.

  \item[H\'eroe] El personaje controlado por cada jugador. Hay m\'as de
    120 h\'eroes disponibles, cada uno con habilidades \'unicas. Los equipos
    eligen sus h\'eroes antes de empezar y no pueden repetir el mismo h\'eroe
    en una partida.

  \item[Draft / Drafting] El proceso de selecci\'on y veto de h\'eroes antes
    de cada partida, equivalente a las rotaciones de jugadores en deportes
    convencionales. Cada equipo elige su estrategia seleccionando h\'eroes
    y bloqueando los del rival.

  \item[Roshan] El enemigo neutral m\'as poderoso del mapa, ubicado en una
    cueva al centro. Matarlo otorga al equipo el \textbf{Aegis of the Immortal}
    (que permite resucitar gratis una vez) y una gran recompensa de oro y
    experiencia. Es el objetivo secundario m\'as relevante del juego.

  \item[Torres (Towers)] Estructuras defensivas que protegen el camino hacia el
    Ancient. Hay varias torres por cada uno de los tres carriles del mapa y
    el equipo atacante debe destruirlas en orden para llegar a la base rival.

  \item[Ward (Observer / Sentry)] Dispositivos de visi\'on que se colocan en
    puntos estrat\'egicos del mapa. La \textbf{Observer ward} otorga visi\'on
    pasiva de un \'area; la \textbf{Sentry ward} revela unidades invisibles y
    wards enemigas ocultas. Controlar la visi\'on del mapa permite planear
    emboscadas y tomar objetivos sin ser sorprendido.

  \item[GPM (Gold Per Minute)] Oro ganado por el equipo por minuto de juego.
    Con el oro se compran \textit{items} (equipamiento) que hacen al h\'eroe
    m\'as poderoso. Un GPM alto significa que el equipo est\'a generando recursos
    de forma eficiente.

  \item[XPM (Experience Per Minute)] Experiencia ganada por el equipo por
    minuto. La experiencia aumenta el nivel del h\'eroe, lo que desbloquea
    habilidades m\'as poderosas. Un XPM alto indica que el equipo crece m\'as
    r\'apido que el rival.

  \item[KDA (Kills / Deaths / Assists)] M\'etrica sint\'etica de rendimiento de
    combate: (Kills + Asistencias) $\div$ Muertes propias. Un KDA de 3.0 significa
    que por cada muerte propia el equipo obtuvo 3 eliminaciones. KDA alto indica
    dominancia en los enfrentamientos directos.

  \item[Farm / Farmear] La acci\'on de matar creeps y neutrales para acumular
    oro y experiencia de forma eficiente, sin necesariamente buscar el combate
    contra los h\'eroes rivales.

  \item[Gank] Emboscada coordinada a uno o m\'as h\'eroes rivales en una posici\'on
    donde no esperan ser atacados. Requiere visi\'on del mapa para ejecutarse
    con seguridad.

  \item[Last Hit / Deny] Mec\'anica exclusiva de Dota~2. El oro de un creep solo
    se obtiene si el jugador da el \textbf{\'ultimo golpe} (last hit). Adem\'as,
    se puede ``denegar'' (deny) un creep aliado para evitar que el rival
    reciba el oro cuando lo mate.

  \item[Creeps / Minions] Tropas autom\'aticas que aparecen cada 30 segundos y
    caminan hacia la base rival. Son la principal fuente de oro y experiencia
    durante los primeros minutos.

  \item[Buyback] Opci\'on de gastar oro para resucitar inmediatamente en lugar de
    esperar. Se usa en momentos cr\'iticos para defender ante un ataque inesperado.

  \item[First Blood] La primera eliminaci\'on de h\'eroe de la partida. Otorga
    una recompensa de oro extra y suele indicar qui\'en lleva la ventaja inicial.

  \item[Meta] Del ingl\'es \textit{Most Effective Tactics Available}: el conjunto
    de estrategias, h\'eroes e \'items considerados \'optimos en un parche
    espec\'ifico. El meta cambia con cada actualizaci\'on del juego.

  \item[Parche (Patch)] Actualizaci\'on del juego publicada por Valve que puede
    cambiar las estad\'isticas de h\'eroes, el mapa o las mec\'anicas, alterando
    radicalmente el meta competitivo.

  \item[Siegeo / Siege] Estrategia de asedio sistem\'atico a las torres enemigas
    sin necesariamente buscar el combate directo. Implica aprovechar la
    ventaja de n\'umero de tropas para presionar estructuras.

  \item[APM (Actions Per Minute)] N\'umero de acciones (clics, comandos) que un
    jugador ejecuta por minuto. Indicador de la velocidad y precisi\'on mec\'anica.

\end{description}

\subsection*{T\'erminos de an\'alisis de datos}

\begin{description}[leftmargin=4.2cm, style=nextline, itemsep=5pt]

  \item[Dataset] Una tabla estructurada de datos. Aqu\'i, cada fila es una partida
    y cada columna es una caracter\'istica medida (GPM, torres, etc.).

  \item[Variable objetivo] La columna que se quiere predecir. Aqu\'i es
    \texttt{radiant\_win}: si gan\'o Radiant (verdadero / falso).

  \item[Feature / Atributo] Una columna del dataset usada como entrada para el
    an\'alisis. Por ejemplo, el GPM o el KDA son features.

  \item[Ventana hist\'orica] Las $N$ partidas previas promediadas para calcular
    las estad\'isticas de cada equipo. Si \texttt{freq\_r = 60}, las estad\'isticas
    de Radiant son el promedio de sus \'ultimas 60 partidas.

  \item[Pick rate] La fracci\'on de partidas (dentro de la ventana hist\'orica) en
    la que un h\'eroe fue elegido. Un pick rate de 0.30 significa que el h\'eroe
    se selecci\'on\'o en el 30\,\% de las partidas recientes del equipo.

  \item[Imputaci\'on] El proceso de rellenar valores faltantes con una estimaci\'on.
    Aqu\'i se us\'o la \textbf{media columnar}: si un equipo no tiene historial,
    se le asigna el promedio global de todos los equipos para esa estad\'istica.

  \item[KNN (K-Nearest Neighbors)] Algoritmo que encuentra las $K$ partidas m\'as
    similares a una consulta midiendo distancias. Sirve para responder
    ``¿cu\'ales partidas del pasado se parecen m\'as a esta?''.

  \item[BallTree] Estructura de datos que organiza los registros en jerarqu\'ia
    para acelerar la b\'usqueda de vecinos cercanos. Sin ella, comparar una
    partida contra 47\,150 ser\'ia lento; con BallTree la b\'usqueda es instant\'anea.

  \item[Correlaci\'on de Pearson] N\'umero entre $-1$ y $+1$ que mide cu\'an
    relacionadas est\'an dos variables linealmente. $+1$: si una sube, la otra
    tambi\'en. $-1$: si una sube, la otra baja. $0$: no hay relaci\'on lineal.

  \item[StandardScaler] Transformaci\'on que lleva todas las variables al mismo
    rango (media 0, desviaci\'on est\'andar 1). Necesaria porque el GPM ($\sim$400)
    y los kills de Roshan ($\sim$0.3) est\'an en escalas muy diferentes; sin
    normalizar, el KNN dar\'ia m\'as peso al GPM solo por ser un n\'umero mayor.

  \item[EDA] \textit{Exploratory Data Analysis} --- An\'alisis Exploratorio de
    Datos: el proceso de examinar el dataset con gr\'aficos y estad\'isticas
    descriptivas antes de construir modelos, para entender su estructura y
    formular hip\'otesis.

  \item[Valor nulo / NaN] Un valor faltante en el dataset. ``NaN'' viene de
    \textit{Not a Number} y significa que esa celda est\'a vac\'ia.

  \item[Outlier] Un valor extremo alejado del resto. Por ejemplo, una partida
    con GPM de 1\,500 cuando la mayor\'ia tiene 400--600 ser\'ia un outlier.

  \item[Distribuci\'on de densidad] Curva que muestra con qu\'e frecuencia aparecen
    los diferentes valores de una variable. Permite comparar grupos de tama\~nos
    distintos en la misma escala, a diferencia del histograma de frecuencias.

\end{description}

\newpage
\section*{Referencias}
\addcontentsline{toc}{section}{Referencias}

\begin{enumerate}
  \item\label{ref:kaggle} Calvo, T. (2021). \textit{Dota 2 Pro Players Matches
        2019--2021} [Dataset]. Kaggle.
        Disponible en: \url{https://www.kaggle.com/datasets/teocalvo/dota2-pro-players-matches-2019-202106/data}
  \item\label{ref:paper} Calvo, T. et al.\ (2021). \textit{Predicting the Outcome
        of Dota 2 Professional Matches}. arXiv preprint arXiv:2106.01782.
        Disponible en: \url{https://arxiv.org/pdf/2106.01782}
  \item\label{ref:sklearn} Pedregosa et al.\ \textit{Scikit-learn: Machine
        Learning in Python}. Journal of Machine Learning Research, 12, 2011.
  \item\label{ref:valve} Valve Corporation. \textit{Dota 2 --- Portal oficial}.
        Disponible en: \url{https://www.dota2.com/}
\end{enumerate}

\end{document}
